\documentclass[a4paper,12pt]{article}

\usepackage[utf8]{inputenc}
\usepackage[T1]{fontenc}
\usepackage[french]{babel}
\usepackage{geometry}
\usepackage{graphicx}
\usepackage{hyperref}
\usepackage{float}
\usepackage{listings}

\geometry{margin=2.5cm}

\title{Pokedle \\ Rapport de projet}
\author{Rémi Soulier \\ BTS SIO SLAM}
\date{\today}

\begin{document}

    \maketitle
    \tableofcontents
    \newpage

    \section{Introduction}

    Dans le cadre de ma formation en BTS SIO option SLAM, j’ai réalisé une application web nommée \textbf{Pokedle}.
    Ce projet étudiant s’inspire du concept du site pokedle.com, qui lui-même reprend le principe du jeu Wordle.

    Le principe du jeu est simple : un Pokémon mystère est choisi aléatoirement et l’utilisateur doit le deviner en un nombre limité d’essais. À chaque tentative, l’application compare les caractéristiques du Pokémon proposé avec celles du Pokémon cible et affiche des indices afin d’aider le joueur.

    Ce projet a été réalisé dans un objectif pédagogique afin de mettre en pratique les compétences du référentiel BTS SIO, notamment :

    \begin{itemize}
        \item la conception d’une application web
        \item la modélisation d’une base de données
        \item l’exploitation d’une API REST
        \item la structuration d’un projet logiciel
    \end{itemize}

    \section{Conception}

    Avant la phase de développement, une phase de conception a été réalisée afin de modéliser le fonctionnement du système.

    \subsection{Diagramme de cas d'utilisation}

    Le diagramme de cas d'utilisation permet d'identifier les interactions entre les utilisateurs et l'application.

    \begin{figure}[H]
        \centering
        \includegraphics[width=12cm]{figures/usecase.png}
        \caption{Diagramme de cas d'utilisation}
    \end{figure}

    Deux types d'utilisateurs sont identifiés :

    \begin{itemize}
        \item le visiteur
        \item le joueur connecté
    \end{itemize}

    Les principales actions possibles sont :

    \begin{itemize}
        \item créer un compte
        \item se connecter
        \item lancer une partie
        \item proposer un Pokémon
        \item consulter le résultat
        \item consulter l’historique des parties
    \end{itemize}

    \subsection{Diagramme de classes}

    Le diagramme de classes permet de représenter la structure de l’application.

    \begin{figure}[H]
        \centering
        \includegraphics[width=12cm]{figures/classdiagram.png}
        \caption{Diagramme de classes}
    \end{figure}

    Les principales classes sont :

    \begin{itemize}
        \item Player : représente un joueur
        \item Game : représente une partie
        \item Pokemon : représente un Pokémon
        \item Attempt : représente une tentative d’un joueur
    \end{itemize}

    Ces classes permettent de modéliser la logique métier du jeu.

    \subsection{Modèle conceptuel de données}

    Le modèle conceptuel de données permet de structurer les informations stockées dans la base de données.

    \begin{figure}[H]
        \centering
        \includegraphics[width=12cm]{figures/mcd.png}
        \caption{Modèle conceptuel de données}
    \end{figure}

    Les entités principales sont :

    \begin{itemize}
        \item Joueur
        \item Partie
        \item Pokémon
        \item Tentative
    \end{itemize}

    Les relations principales sont :

    \begin{itemize}
        \item un joueur peut jouer plusieurs parties
        \item une partie contient plusieurs tentatives
        \item chaque tentative correspond à un Pokémon proposé
    \end{itemize}

    \section{Stack logicielle}

    L'application repose sur plusieurs technologies.

    \subsection{Backend}

    Le backend est développé en PHP en utilisant le framework Laravel.

    Laravel offre plusieurs avantages :

    \begin{itemize}
        \item architecture MVC
        \item gestion simplifiée des routes
        \item ORM Eloquent pour la base de données
        \item gestion de l’authentification
    \end{itemize}

    \subsection{Frontend}

    L’interface utilisateur est développée avec :

    \begin{itemize}
        \item HTML
        \item CSS
        \item Blade (moteur de template Laravel)
        \item Tailwind CSS
    \end{itemize}

    Blade permet d’organiser les vues et de réutiliser des composants communs comme le layout principal.

    \subsection{Base de données}

    La base de données utilisée est MySQL.

    Elle est gérée dans Laravel grâce à :

    \begin{itemize}
        \item les migrations
        \item les seeders
        \item l’ORM Eloquent
    \end{itemize}

    Les migrations permettent de versionner la structure de la base de données.

    \subsection{API externe}

    Les données des Pokémon sont récupérées via l’API publique PokéAPI.

    Cette API permet de récupérer :

    \begin{itemize}
        \item les caractéristiques des Pokémon
        \item leurs types
        \item leurs évolutions
        \item leurs statistiques
    \end{itemize}

    Ces données sont ensuite stockées dans la base afin d’optimiser les performances de l’application.

    \section{Présentation des sections de l'application}

    \subsection{Page d'accueil}

    La page d'accueil présente le concept du jeu et permet à l'utilisateur :

    \begin{itemize}
        \item de se connecter
        \item de créer un compte
        \item de jouer en mode invité
    \end{itemize}

    \subsection{Authentification}

    L'application propose un système d'authentification permettant :

    \begin{itemize}
        \item la création d’un compte
        \item la connexion
        \item la gestion des sessions utilisateur
    \end{itemize}

    Les mots de passe sont sécurisés grâce à un système de hachage.

    \subsection{Interface de jeu}

    L’interface de jeu permet à l’utilisateur de proposer un Pokémon.

    À chaque tentative :

    \begin{itemize}
        \item l'application récupère les informations du Pokémon proposé
        \item elle compare ces informations avec celles du Pokémon cible
        \item elle affiche les résultats sous forme d'indices
    \end{itemize}

    Les indices permettent d’indiquer si les caractéristiques sont correctes ou non.

    \subsection{Gestion des parties}

    Chaque partie est enregistrée dans la base de données.

    Les informations enregistrées sont :

    \begin{itemize}
        \item le joueur
        \item le Pokémon à deviner
        \item le nombre d'essais
        \item la date de la partie
    \end{itemize}

    \section{Points d'amélioration}

    Plusieurs améliorations peuvent être envisagées pour faire évoluer le projet.

    \subsection{Mode de jeu quotidien}

    Un mode quotidien pourrait être ajouté afin de proposer le même Pokémon à tous les joueurs chaque jour.

    \subsection{Classement des joueurs}

    Un système de classement pourrait être mis en place afin de comparer les performances des joueurs.

    \subsection{Amélioration de l'interface}

    L’interface graphique pourrait être améliorée afin de proposer une expérience utilisateur plus immersive.

    \subsection{Statistiques}

    Il serait possible d'ajouter des statistiques telles que :

    \begin{itemize}
        \item le nombre de parties jouées
        \item le taux de réussite
        \item le nombre moyen d’essais
    \end{itemize}

    \subsection{Déploiement}

    L’application pourrait être déployée sur un serveur afin d’être accessible publiquement.

    \section{Conclusion}

    Le projet Pokedle m’a permis de concevoir et développer une application web complète en utilisant le framework Laravel.

    Ce projet m’a permis de renforcer mes compétences en :

    \begin{itemize}
        \item développement web
        \item conception de bases de données
        \item utilisation d’API REST
        \item organisation d’un projet informatique
    \end{itemize}

    Il constitue une application fonctionnelle illustrant les compétences acquises durant ma formation en BTS SIO option SLAM.

\end{document}
